\begin{corollary}[7.1.13]
\label{IV.2.s7.1.13-fr}
Sean $Y$ un preesquema localmente noetheriano
irreducible de dimensión \textup{1}, $X$ un preesquema
irreducible, $f:X\to Y$ un morfismo dominante de
tipo finito, y $\xi$, $\eta$ los puntos genéricos
de $X$ e $Y$, respectivamente. Entonces, para
todo $y\in f(X)$, las dimensiones de las componentes
irreducibles de $f^{-1}(y)$ son todas iguales a
$\deg.\operatorname{tr}_{\mathbf{k}(\eta)}\boldsymbol{k}(\xi)$.
\end{corollary}

\begin{proof}
Por hipótesis, $f^{-1}(\eta)$ es irreducible, tiene
por punto genérico a $\xi$ y su dimensión es
$\deg.\operatorname{tr}_{\mathbf{k}(\eta)}\boldsymbol{k}(\xi)$ (5.2.1).
Como $Y$ es irreducible y de dimensión 1, se tiene
$\dim(\mathscr O_y)=1$ para todo $y\neq\eta$, y
$\mathscr O_y$ es universalmente catenario en virtud
de (7.1.12, (ii)). Si $y\in f(X)$ y $z$ es un punto
genérico de una componente irreducible $Z$ de
$f^{-1}(y)$, se tiene, por (5.6.5),
\[
\dim(Z)=1+\dim(f^{-1}(\eta))\setmin\dim(\mathscr O_{X,z}).
\]
Pero $\dim(\mathscr O_{X,z})\leqslant\dim(\mathscr O_y)$
(\textbf{0}, 16.3.9). Por otra parte, como $z$ no es el punto
genérico de $X$, $\dim(\mathscr O_{X,z})>0$; por tanto,
$\dim(\mathscr O_{X,z})=1$ y
$\dim(Z)=\dim(f^{-1}(\eta))$.
\end{proof}
